\section{Details: Infinite products (part 2) -- Finite product composition rule}

This chapter contains the detailed proof of the finite product composition
rule for formal power series (Lemma \ref{lem.fps.subs.rule-infprod-fin}),
which states that substitution (composition) distributes over finite products.
This lemma is used in the proof of Proposition \ref{prop.fps.subs.rule-infprod}
(the infinite product composition rule), whose detailed proof is also given here.

\begin{lemma}
\label{lem.fps.subs.rule-infprod-fin}
\lean{PowerSeries.comp_prod_finite}
\leantarget
\leanok
Let $I$ be a \textbf{finite} set. If
$\left(  f_{i}\right)  _{i\in I}\in K\left[  \left[  x\right]  \right]  ^{I}$
is a family of FPSs, and if $g\in K\left[  \left[  x\right]  \right]  $ is an
FPS satisfying $\left[  x^{0}\right]  g=0$, then $\left(  \prod_{i\in I}%
f_{i}\right)  \circ g=\prod_{i\in I}\left(  f_{i}\circ g\right)  $.
\end{lemma}

\begin{proof}
\leanok
This follows by a straightforward induction on $\left\vert I\right\vert $.

\textbf{Base case} ($\left\vert I\right\vert =0$):
The empty product is $\underline{1}$ (the constant FPS with value $1$),
and we have $\underline{1}\circ g=\underline{1}$ for any
$g\in K\left[  \left[  x\right]  \right]  $.
Hence both sides equal $\underline{1}$.

\textbf{Induction step:}
Assume the result holds for all sets of size $\left\vert I\right\vert - 1$.
Pick some $a\in I$ and write
$\prod_{i\in I}f_{i}=f_{a}\cdot\prod_{i\in I\setminus\left\{  a\right\}  }f_{i}$.
By Proposition \ref{prop.fps.subs.rules} \textbf{(b)}
(the substitution rule for products:
$\left(  u\cdot v\right)  \circ g=\left(  u\circ g\right)  \cdot\left(
v\circ g\right)  $), we have
\[
\left(  f_{a}\cdot\prod_{i\in I\setminus\left\{  a\right\}  }f_{i}\right)
\circ g=\left(  f_{a}\circ g\right)  \cdot\left(  \prod_{i\in
I\setminus\left\{  a\right\}  }f_{i}\right)  \circ g.
\]
By the induction hypothesis (applied to the set
$I\setminus\left\{  a\right\}  $, which has size
$\left\vert I\right\vert -1$),
\[
\left(  \prod_{i\in I\setminus\left\{  a\right\}  }f_{i}\right)  \circ
g=\prod_{i\in I\setminus\left\{  a\right\}  }\left(  f_{i}\circ g\right)  .
\]
Combining, we obtain
\[
\left(  \prod_{i\in I}f_{i}\right)  \circ g=\left(  f_{a}\circ g\right)
\cdot\prod_{i\in I\setminus\left\{  a\right\}  }\left(  f_{i}\circ g\right)
=\prod_{i\in I}\left(  f_{i}\circ g\right)  .
\]
\end{proof}

\subsection{Detailed proof of Proposition~\ref{prop.fps.subs.rule-infprod}}

\begin{proof}[Detailed proof of Proposition \ref{prop.fps.subs.rule-infprod}.]
\proves{prop.fps.subs.rule-infprod}
\leanok
Let $\left(
f_{i}\right)  _{i\in I}\in K\left[  \left[  x\right]  \right]  ^{I}$ be a
multipliable family of FPSs. Let $g\in K\left[  \left[  x\right]  \right]  $
be an FPS satisfying $\left[  x^{0}\right]  g=0$.

We shall first show the following auxiliary claim:

\textit{Claim 1:} Let $M$ be an $x^{n}$-approximator for $\left(
f_{i}\right)  _{i\in I}$. Then, the set $M$ determines the $x^{n}$-coefficient
in the product of $\left(  f_{i}\circ g\right)  _{i\in I}$.

[\textit{Proof of Claim 1:} The set $M$ is an $x^{n}$-approximator for
$\left(  f_{i}\right)  _{i\in I}$. In other words, $M$ is a finite subset of
$I$ that determines the first $n+1$ coefficients in the product of $\left(
f_{i}\right)  _{i\in I}$ (by the definition of ``$x^{n}$-approximator'').

Let $J$ be a finite subset of $I$ satisfying $M\subseteq J\subseteq I$. Then,
Lemma \ref{lem.fps.subs.rule-infprod-fin} (applied to $J$ instead of $I$)
yields%
\begin{equation}
\left(  \prod_{i\in J}f_{i}\right)  \circ g=\prod_{i\in J}\left(  f_{i}\circ
g\right)  . \label{pf.prop.fps.subs.rule-infprod.3}%
\end{equation}
Also, Lemma \ref{lem.fps.subs.rule-infprod-fin} (applied to $M$ instead of
$I$) yields%
\begin{equation}
\left(  \prod_{i\in M}f_{i}\right)  \circ g=\prod_{i\in M}\left(  f_{i}\circ
g\right)  . \label{pf.prop.fps.subs.rule-infprod.4}%
\end{equation}
However, Proposition \ref{prop.fps.infprod-approx-xneq} \textbf{(a)} (applied
to $\mathbf{a}_{i}=f_{i}$) yields%
\[
\prod_{i\in J}f_{i}\overset{x^{n}}{\equiv}\prod_{i\in M}f_{i}.
\]
We also have $g\overset{x^{n}}{\equiv}g$ (since the relation $\overset{x^{n}%
}{\equiv}$ is an equivalence relation). Hence, Proposition
\ref{prop.fps.xneq.comp} (applied to $a=\prod_{i\in J}f_{i}$ and
$b=\prod_{i\in M}f_{i}$ and $c=g$ and $d=g$) yields%
\[
\left(  \prod_{i\in J}f_{i}\right)  \circ g\overset{x^{n}}{\equiv}\left(
\prod_{i\in M}f_{i}\right)  \circ g.
\]
In view of (\ref{pf.prop.fps.subs.rule-infprod.3}) and
(\ref{pf.prop.fps.subs.rule-infprod.4}), this rewrites as%
\[
\prod_{i\in J}\left(  f_{i}\circ g\right)  \overset{x^{n}}{\equiv}\prod_{i\in
M}\left(  f_{i}\circ g\right)  .
\]
In other words,
\[
\text{each }m\in\left\{  0,1,\ldots,n\right\}  \text{ satisfies }\left[
x^{m}\right]  \left(  \prod_{i\in J}\left(  f_{i}\circ g\right)  \right)
=\left[  x^{m}\right]  \left(  \prod_{i\in M}\left(  f_{i}\circ g\right)
\right)
\]
(by the definition of the relation $\overset{x^{n}}{\equiv}$). Applying this
to $m=n$, we obtain%
\[
\left[  x^{n}\right]  \left(  \prod_{i\in J}\left(  f_{i}\circ g\right)
\right)  =\left[  x^{n}\right]  \left(  \prod_{i\in M}\left(  f_{i}\circ
g\right)  \right)  .
\]

Forget that we fixed $J$. We thus have shown that every finite subset $J$ of
$I$ satisfying $M\subseteq J\subseteq I$ satisfies%
\[
\left[  x^{n}\right]  \left(  \prod_{i\in J}\left(  f_{i}\circ g\right)
\right)  =\left[  x^{n}\right]  \left(  \prod_{i\in M}\left(  f_{i}\circ
g\right)  \right)  .
\]
In other words, the set $M$ determines the $x^{n}$-coefficient in the product
of $\left(  f_{i}\circ g\right)  _{i\in I}$ (by the definition of what it
means to ``determine the $x^{n}$-coefficient in the product of
$\left(  f_{i}\circ g\right)  _{i\in I}$''). This proves Claim
1.] \medskip

Now, let $n\in\mathbb{N}$. Lemma \ref{lem.fps.mulable.approx} (applied to
$\mathbf{a}_{i}=f_{i}$) shows that there exists an $x^{n}$-approximator for
$\left(  f_{i}\right)  _{i\in I}$. Consider this $x^{n}$-approximator for
$\left(  f_{i}\right)  _{i\in I}$, and denote it by $M$. Thus, $M$ is an
$x^{n}$-approximator for $\left(  f_{i}\right)  _{i\in I}$; in other words,
$M$ is a finite subset of $I$ that determines the first $n+1$ coefficients in
the product of $\left(  f_{i}\right)  _{i\in I}$ (by the definition of
``$x^{n}$-approximator''). Claim 1 shows that
the set $M$ determines the $x^{n}$-coefficient in the product of $\left(
f_{i}\circ g\right)  _{i\in I}$. Hence, there is a finite subset of $I$ that
determines the $x^{n}$-coefficient in the product of $\left(  f_{i}\circ
g\right)  _{i\in I}$ (namely, $M$). In other words, the $x^{n}$-coefficient in
the product of $\left(  f_{i}\circ g\right)  _{i\in I}$ is finitely determined.

Forget that we fixed $n$. We thus have shown that for each $n\in\mathbb{N}$,
the $x^{n}$-coefficient in the product of $\left(  f_{i}\circ g\right)  _{i\in
I}$ is finitely determined. In other words, each coefficient in the product of
$\left(  f_{i}\circ g\right)  _{i\in I}$ is finitely determined. In other
words, the family $\left(  f_{i}\circ g\right)  _{i\in I}$ is multipliable (by
the definition of ``multipliable''). \medskip

It remains to prove that $\left(  \prod_{i\in I}f_{i}\right)  \circ
g=\prod_{i\in I}\left(  f_{i}\circ g\right)  $.

In order to do so, we again fix $n\in\mathbb{N}$. Lemma
\ref{lem.fps.mulable.approx} (applied to $\mathbf{a}_{i}=f_{i}$) shows that
there exists an $x^{n}$-approximator for $\left(  f_{i}\right)  _{i\in I}$.
Consider this $x^{n}$-approximator for $\left(  f_{i}\right)  _{i\in I}$, and
denote it by $M$. Thus, $M$ is an $x^{n}$-approximator for $\left(
f_{i}\right)  _{i\in I}$; in other words, $M$ is a finite subset of $I$ that
determines the first $n+1$ coefficients in the product of $\left(
f_{i}\right)  _{i\in I}$ (by the definition of ``$x^{n}%
$-approximator''). Moreover, Proposition
\ref{prop.fps.infprod-approx-xneq} \textbf{(b)} (applied to $\mathbf{a}%
_{i}=f_{i}$) yields%
\begin{equation}
\prod_{i\in I}f_{i}\overset{x^{n}}{\equiv}\prod_{i\in M}f_{i}.
\label{pf.prop.fps.subs.rule-infprod.1}%
\end{equation}
We also have $g\overset{x^{n}}{\equiv}g$ (since the relation $\overset{x^{n}%
}{\equiv}$ is an equivalence relation). Hence, Proposition
\ref{prop.fps.xneq.comp} (applied to $a=\prod_{i\in I}f_{i}$ and
$b=\prod_{i\in M}f_{i}$ and $c=g$ and $d=g$) yields%
\begin{equation}
\left(  \prod_{i\in I}f_{i}\right)  \circ g\overset{x^{n}}{\equiv}\left(
\prod_{i\in M}f_{i}\right)  \circ g. \label{pf.prop.fps.subs.rule-infprod.2}%
\end{equation}
However, Lemma \ref{lem.fps.subs.rule-infprod-fin} (applied to $M$ instead of
$I$) yields%
\[
\left(  \prod_{i\in M}f_{i}\right)  \circ g=\prod_{i\in M}\left(  f_{i}\circ
g\right)  .
\]
In view of this, we can rewrite (\ref{pf.prop.fps.subs.rule-infprod.2}) as
\[
\left(  \prod_{i\in I}f_{i}\right)  \circ g\overset{x^{n}}{\equiv}\prod_{i\in
M}\left(  f_{i}\circ g\right)  .
\]
In other words,
\[
\text{each }m\in\left\{  0,1,\ldots,n\right\}  \text{ satisfies }\left[
x^{m}\right]  \left(  \left(  \prod_{i\in I}f_{i}\right)  \circ g\right)
=\left[  x^{m}\right]  \left(  \prod_{i\in M}\left(  f_{i}\circ g\right)
\right)
\]
(by the definition of the relation $\overset{x^{n}}{\equiv}$). Applying this
to $m=n$, we obtain%
\begin{equation}
\left[  x^{n}\right]  \left(  \left(  \prod_{i\in I}f_{i}\right)  \circ
g\right)  =\left[  x^{n}\right]  \left(  \prod_{i\in M}\left(  f_{i}\circ
g\right)  \right)  . \label{pf.prop.fps.subs.rule-infprod.6}%
\end{equation}

However, Claim 1 shows that the set $M$ determines the $x^{n}$-coefficient in
the product of $\left(  f_{i}\circ g\right)  _{i\in I}$. Hence, the definition
of the infinite product $\prod_{i\in I}\left(  f_{i}\circ g\right)  $
(specifically, Definition \ref{def.fps.multipliable} \textbf{(b)}) yields%
\[
\left[  x^{n}\right]  \left(  \prod_{i\in I}\left(  f_{i}\circ g\right)
\right)  =\left[  x^{n}\right]  \left(  \prod_{i\in M}\left(  f_{i}\circ
g\right)  \right)  .
\]
Comparing this with (\ref{pf.prop.fps.subs.rule-infprod.6}), we obtain%
\[
\left[  x^{n}\right]  \left(  \left(  \prod_{i\in I}f_{i}\right)  \circ
g\right)  =\left[  x^{n}\right]  \left(  \prod_{i\in I}\left(  f_{i}\circ
g\right)  \right)  .
\]

Forget that we fixed $n$. We thus have shown that each $n\in\mathbb{N}$
satisfies%
\[
\left[  x^{n}\right]  \left(  \left(  \prod_{i\in I}f_{i}\right)  \circ
g\right)  =\left[  x^{n}\right]  \left(  \prod_{i\in I}\left(  f_{i}\circ
g\right)  \right)  .
\]
In other words, the FPSs $\left(  \prod_{i\in I}f_{i}\right)  \circ g$ and
$\prod_{i\in I}\left(  f_{i}\circ g\right)  $ agree in all their coefficients.
Hence, $\left(  \prod_{i\in I}f_{i}\right)  \circ g=\prod_{i\in I}\left(
f_{i}\circ g\right)  $. This completes our proof of Proposition
\ref{prop.fps.subs.rule-infprod}.
\end{proof}
